\subsection{Application: invariants of a group of automorphisms of an algebra}

Given a ring K, a K-algebra A and a group G, we shall say that G operates on A if: (1) the set A has group of operators G (Algebra, Chapter I, § 7, no. 2); (2) for all \(\sigma \in G\) , the mapping \(x \mapsto \sigma \cdot x\) is an endomorphism of the K-algebra A (and therefore an automorphism since it is bijective (loc. cit.)). We shall denote by \(A^{G}\) , the set of elements of A which are invariant under G; clearly it is a sub-K-algebra of A.

We shall say that \(\mathcal{G}\) is a locally finite group of operators on A if every orbit of \(\mathcal{G}\) in A (Algebra, Chapter I, Corrections to Fascicule IV) is finite.

PROPOSITION 22. Let A be a (commutative) K-algebra and G a locally finite group of operators on A. Then A is integral over the subalgebra \(\mathbf{A}^{\mathcal{G}}\) .

For all \(x \in \mathbf{A}\) , let \(x_{i} (1 \leqslant i \leqslant n)\) be the distinct elements of the orbit of \(x\) under \(\mathcal{G}\) ; for all \(\sigma \in \mathcal{G}\) , there exists a permutation \(\pi_{\sigma}\) of the set \(\{1, 2, \ldots, n\}\) such that \(\sigma. x_{i} = x_{\pi_{\sigma}(i)}\) for \(1 \leqslant i \leqslant n\) ; therefore the elementary symmetric functions of the \(x_{i}\) are elements of \(\mathbf{A}\) which are invariant under \(\mathcal{G}\) , in other

words elements of \(\mathbf{A}^{\mathcal{G}}\) . As \(x\) is a root of the monic polynomial \(\prod_{i=1}^{n} (\mathbf{X} - x_i)\) and the coefficients of this polynomial belong to \(\mathbf{A}^{\mathcal{G}}\) , \(x\) is integral over \(\mathbf{A}^{\mathcal{G}}\) .

THEOREM 2. Let A be a finitely generated K-algebra and G a locally finite group of operators on A. Then A is a finitely generated \(\mathbf{A}^{\mathcal{G}}\)-module; if further K is Noetherian, \(A^{G}\) is a finitely generated K-algebra.

Let \((a_{j})_{1\leqslant j\leqslant m}\) be a system of generators of the K-algebra A; as a fortiori \(\mathbf{A} = \mathbf{A}^{\mathcal{G}}[a_1,\ldots ,a_m]\) and the \(a_{j}\) are integral over \(\mathbf{A}^{\mathcal{G}}\) by Proposition 22, the first assertion follows from no. 1, Proposition 4. The second is a consequence of the following lemma:

LEMMA 5. Let K be a Noetherian ring, B a finitely generated K-algebra and C a sub-K-algebra of B such that B is integral over C. Then C is a finitely generated K-algebra.

Let \((x_{i})_{1\leqslant i\leqslant n}\) be a finite system of generators of the K-algebra B. For all \(i\) , there exists by hypothesis a monic polynomial \(\mathbf{P}_i\in \mathbf{C}[\mathbf{X}]\) such that \(\mathrm{P_i}(x_i) = 0\) . Let \(\mathbf{C}'\) be the sub-K-algebra of C generated by the coefficients of the \(\mathbf{P}_i\) ( \(1\leqslant i\leqslant n\) ); clearly the \(x_{i}\) are integral over \(\mathbf{C}'\) and \(\mathbf{B} = \mathbf{C}'[x_1,\dots,x_n]\) ; hence B is a finitely generated \(\mathbf{C}'\) -module (no. 1, Proposition 4). On the other hand \(\mathbf{C}'\) is a Noetherian ring (Chapter III, §2, no. 10, Corollary 3 to Theorem 2); hence C is a finitely generated \(\mathbf{C}'\) -module, which proves that C is a finitely generated K-algebra.

Remark. The set of \(\sigma \in \mathcal{G}\) such that \(\sigma a_{j} = a_{j}\) for \(1 \leqslant j \leqslant m\) obviously leaves invariant every element of A. The normal subgroup \(\mathcal{H}\) of \(\mathcal{G}\) leaving invariant every element of A is therefore of finite index in \(\mathcal{G}\) and A may be considered as having a finite group of operators \(\mathcal{G}/\mathcal{H}\) ; obviously \(\mathrm{A}^{\mathcal{G}/\mathcal{H}} = \mathrm{A}^{\mathcal{G}}\) .

Let S be a multiplicative subset of a ring A and G a group operating on A and for which S is stable; then, for all \(\sigma\in G\) , there exists a unique endomorphism \(z\mapsto\sigma.z\) of the ring \(S^{-1}A\) such that \(\sigma.(a/1)=(\sigma.a)/1\) for all \(a\in A\) ; it is given by the formula \(\sigma.(a/s)=(\sigma.a)/( \sigma.s)\) for \(a\in A\) and \(s\in S\) (Chapter II, §2, no. 1, Proposition 2); if \(\tau\) is another element of G, clearly \(\sigma.(\tau.z)=(\sigma\tau).z\) for all \(z\in S^{-1}A\) and hence the group G operates on the ring \(S^{-1}A\) .

PROPOSITION 23. Let A be a K-algebra, \(\mathcal{G}\) a locally finite group of operators on A, S a multiplicative subset of A stable under \(\mathcal{G}\) and \(\mathrm{S}^{\mathcal{G}}\) the set \(\mathrm{S} \cap \mathrm{A}^{\mathcal{G}}\) . Then the canonical mapping of \((\mathrm{S}^{\mathcal{G}})^{-1}\mathrm{A}\) to \(\mathrm{S}^{-1}\mathrm{A}\) (Chapter II, § 2, no. 1, Corollary 2 to Proposition 2) is an isomorphism which maps \((\mathrm{S}^{\mathcal{G}})^{-1}\mathrm{A}^{\mathcal{G}}\) to \((\mathrm{S}^{-1}\mathrm{A})^{\mathcal{G}}\) .

For all \(s \in S\) , let \(s, s_1, \ldots, s_q\) be the distinct elements of the orbit of \(s\) under \(\mathcal{G}\) ; as \(ss_1 \ldots s_q \in S^{\mathcal{G}}\) , the first assertion follows from Chapter II, § 2, no. 3, Proposition 8. Identifying canonically \((S^{\mathcal{G}})^{-1}A\) with \(S^{-1}A\) , clearly every element of \((S^{\mathcal{G}})^{-1}A^{\mathcal{G}}\) is invariant under \(\mathcal{G}\) . Conversely, let \(a/t\) be an element of \((S^{\mathcal{G}})^{-1}A\) which is invariant under \(\mathcal{G}\) ( \(a \in A\) , \(t \in S^{\mathcal{G}}\) ); if \(a_j\) ( \(1 \leqslant j \leqslant m\) ) are the distinct elements of the orbit of \(a\) under \(\mathcal{G}\) , then \(a_j/t = a/t\) for \(1 \leqslant j \leqslant m\) and therefore there exists \(s \in S^{\mathcal{G}}\) such that \(s(a_j - a) = 0\) for \(1 \leqslant j \leqslant m\) ; in other words, \(sa\) is invariant under \(\mathcal{G}\) and, as \(a/t = (sa)/(st)\) , certainly \(a/t \in (S^{\mathcal{G}})^{-1}A^{\mathcal{G}}\) .

COROLLARY. Let A be an integral domain, K its field of fractions and G a locally finite group of operators on A. Then G operates on K and \(K^{G}\) is the field of fractions of \(A^{G}\) .

A - \{0\} is stable under G.

